\documentclass[12pt,a4paper]{article}
\usepackage[utf8]{inputenc}
\usepackage[spanish]{babel}
\usepackage{geometry}
\usepackage{booktabs}
\usepackage{enumitem}
\usepackage{hyperref}
\geometry{margin=2.5cm}

\title{Actividad de laboratorio: introducción práctica a sistemas P2P\\
\large Laboratorio TIC --- DIINF}
\author{Curso: Sistemas Distribuidos y Paralelos\\
Departamento de Ingeniería Informática --- USACH}
\date{10 de julio de 2026}

\begin{document}

\maketitle

\section{Objetivo de aprendizaje}

En una hora, cada grupo montará una mini-red P2P entre sus nodos, repartirá un archivo en piezas, implementará una estrategia de descarga y medirá tiempos y mensajes ante fallas. El foco está en \textbf{experimentar y medir}, no solo en revisar teoría.

\section{Rol del docente}

El docente orienta con \textbf{teoría} y responde dudas si las solicitan. \textbf{No} prepara archivos ni configura la red: eso lo hace cada grupo.

\section{Repositorio y fork}

Código base y este enunciado están en el repositorio público del curso:

\begin{center}
\url{https://github.com/miguelcarcamov/lab_experiences}
\end{center}

Cada grupo debe:

\begin{enumerate}
    \item Hacer \textbf{fork} del repositorio a su cuenta de GitHub.
    \item Clonar \textbf{su fork} en al menos una máquina del grupo:
    \begin{verbatim}
git clone git@github.com:<usuario>/lab_experiences.git
cd lab_experiences/p2p-intro
    \end{verbatim}
    \item Implementar \texttt{select\_next\_piece} en \texttt{node.py} y commitear en \texttt{master} de su fork.
    \item No subir artefactos generados en el lab (\texttt{demo.bin}, \texttt{payload/}, \texttt{piezas/}); el \texttt{.gitignore} ya los excluye.
\end{enumerate}

\section{Organización}

\begin{itemize}
    \item \textbf{Lugar:} Laboratorio TIC, DIINF.
    \item \textbf{Duración:} 60 minutos.
    \item \textbf{Grupos:} 3 equipos de 5 estudiantes (los mismos del laboratorio del curso).
    \item \textbf{Modalidad:} cada grupo usa varios nodos del lab (recomendado: al menos 3 por grupo). Deben conectarlos en la misma red local (switch) y anotar la IP de cada máquina.
\end{itemize}

\section{Materiales y prerrequisitos}

\begin{itemize}
    \item Nodos del laboratorio TIC y switch para interconectarlos.
    \item Python~3 (solo librería estándar).
    \item Carpeta \texttt{p2p-intro/} del fork (archivos principales):
    \begin{itemize}
        \item \texttt{node.py} --- servidor, protocolo y \texttt{fetch} listos; \textbf{TODO:} \texttt{select\_next\_piece(...)}.
        \item \texttt{discover.py} --- descubrimiento TCP, listo para usar.
        \item \texttt{prepare\_payload.py} --- generan y trocean el archivo de prueba del grupo.
        \item \texttt{activity\_1.pdf} --- este enunciado.
    \end{itemize}
    \item Papel para dibujar el grafo de conectividad.
\end{itemize}

\subsection{Preparación del archivo (lo hace el grupo)}

\begin{enumerate}
    \item Generar un archivo de prueba: \texttt{python3 prepare\_payload.py --gen-demo} (crea \texttt{demo.bin}, $\sim$48\,KB).
    \item Trocearlo y repartir piezas entre los nodos del grupo, por ejemplo con 4 nodos:
    \begin{verbatim}
python3 prepare_payload.py --input demo.bin --chunk-size 4096 --nodes 4
    \end{verbatim}
    Esto produce \texttt{manifest.json} y carpetas \texttt{payload/nodes/node01/}, \ldots con subconjuntos de \texttt{chunk\_NNN.bin}.
    \item Copiar en cada PC su carpeta como \texttt{./piezas/} (cada nodo arranca con \textbf{parte} del archivo, nadie con el archivo completo).
    \item Crear \texttt{seeds.txt} y \texttt{peers.txt} con las IP reales del grupo (\texttt{host:8800:node-id} por línea).
\end{enumerate}

La reconstrucción es correcta si \texttt{recuperado.bin} tiene el mismo SHA-256 que \texttt{manifest.json}.

\subsection{Por qué el ejercicio es \texttt{select\_next\_piece}}

La configuración de red ya está resuelta en \texttt{node.py}. Lo pedagógico es \textbf{decidir a qué peer pedir qué pieza} (secuencial, \emph{rarest-first}, aleatorio, etc.): eso cambia mensajes, tiempos y tolerancia a fallas.

\section{Desarrollo de la actividad}

\begin{center}
\begin{tabular}{@{}clp{8.5cm}@{}}
\toprule
\textbf{Tiempo} & \textbf{Fase} & \textbf{Qué hace cada grupo} \\
\midrule
0--10 min &
Fork + red + descubrimiento &
Fork y clone del repo; conectar nodos; \texttt{prepare\_payload.py}; copiar \texttt{piezas/}; \texttt{node.py serve}; \texttt{seeds.txt}/\texttt{peers.txt}; \texttt{discover.py} y grafo. \\
\addlinespace
10--40 min &
Implementación + difusión &
\textbf{10--15 min:} completar \texttt{select\_next\_piece} y commit en su fork. \textbf{15--40 min:} \texttt{node.py fetch ...}; anotar tiempos y mensajes. \\
\addlinespace
40--50 min &
Inyección de falla &
\textbf{Otro grupo} detiene 1--2 nodos \texttt{serve} sin aviso. Reintentar \texttt{fetch}; anotar detección y recuperación. \\
\addlinespace
50--60 min &
Cierre con datos &
Plenario: estrategia, tiempos, mensajes, grafo, efecto de la falla. URL del fork si el docente lo pide. \\
\bottomrule
\end{tabular}
\end{center}

\subsection{Indicaciones operativas}

\begin{itemize}
    \item Repartan roles (red, payload, implementación, Git, medición, reporte).
    \item Solo modifiquen \texttt{select\_next\_piece} en \texttt{node.py} salvo acuerdo documentado del grupo.
    \item Si \texttt{discover.py} no ve peers, revisen IP, firewall y que \texttt{serve} esté activo.
    \item Documenten la estrategia usada en la medición final.
\end{itemize}

\section{Preguntas de reflexión (cierre grupal)}

\begin{enumerate}
    \item ¿Qué estrategia implementaron en \texttt{select\_next\_piece}? ¿Por qué?
    \item ¿Cómo detectaron un nodo caído?
    \item ¿Cuántos mensajes y cuánto tiempo antes y después de la falla?
    \item Si cayeran la mitad de los nodos, ¿su estrategia seguiría funcionando? ¿Qué faltaría (DHT, réplicas, \emph{rarest-first} global)?
\end{enumerate}

\section{Extensión opcional}

Quienes terminen antes pueden probar \emph{rarest-first} o esbozar cómo localizarían piezas sin \texttt{peers.txt} central.

\vspace{0.5cm}
\noindent\textbf{Entregable mínimo (si aplica):} URL del fork, nodos e IPs, estrategia, tiempos y mensajes, reacción ante la falla, grafo de conectividad.

\end{document}
